%!TEX program = xelatex
\documentclass{beamer}

\usepackage{ctex}
\usepackage[T1]{fontenc}
\usepackage{graphicx}
\usepackage{amssymb,amsmath,amsthm}
\usepackage{array,booktabs,multicol,multirow,tabularx}
\usepackage{listings}
\usepackage[dvipsnames,svgnames]{xcolor}
\usepackage{tcolorbox}
\usepackage{tikz}
\usepackage{hyperref}
\usefonttheme{professionalfonts}
\usepackage{College}
\setmonofont{Fira Code}

\author{ECUSTPC 2026 Fall 裁判组}
\title{2026 年华东理工大学程序设计竞赛（秋季）\\试题分析}
\institute{ECUST}
\date{2026 年秋季}

\definecolor{deepblue}{rgb}{0,0,0.5}
\definecolor{deepred}{rgb}{0.6,0,0}
\definecolor{deepgreen}{rgb}{0,0.5,0}
\definecolor{halfgray}{gray}{0.55}
\newcommand{\contestbox}[2]{%
  \newtcbox{#1}{on line,arc=0pt,outer arc=1pt,colback=#2,colframe=#2,
  boxsep=0pt,left=3pt,right=3pt,top=3pt,bottom=3pt,boxrule=0pt,
  bottomrule=1pt,toprule=1pt,leftrule=1pt,rightrule=1pt}}
\newcommand{\contestboxwhite}[2]{%
  \newtcbox{#1}{on line,arc=0pt,outer arc=1pt,colback=#2,colframe=#2,
  coltext=white,boxsep=0pt,left=3pt,right=3pt,top=3pt,bottom=3pt,
  boxrule=0pt,bottomrule=1pt,toprule=1pt,leftrule=1pt,rightrule=1pt}}

\definecolor{Line1}{HTML}{FFFF00}
\definecolor{Line2}{HTML}{0066FF}
\definecolor{Line3}{HTML}{9966FF}
\definecolor{Line4}{HTML}{33CC33}
\definecolor{Line5}{HTML}{FF5050}
\definecolor{Line6}{HTML}{99CCFF}
\definecolor{Line7}{HTML}{CCCCCC}
\definecolor{Line8}{HTML}{333333}
\definecolor{Line9}{HTML}{FFCC66}
\definecolor{Line10}{HTML}{CC00CC}
\definecolor{Line11}{HTML}{CCFF66}
\definecolor{Line12}{HTML}{3DC3B6}
\definecolor{Line13}{HTML}{F4C2C2}
\contestbox{\myboxi}{Line1}       \contestboxwhite{\myboxii}{Line2}
\contestboxwhite{\myboxiii}{Line3} \contestbox{\myboxiv}{Line4}
\contestboxwhite{\myboxv}{Line5} \contestbox{\myboxvi}{Line6}
\contestbox{\myboxvii}{Line7}    \contestboxwhite{\myboxviii}{Line8}
\contestbox{\myboxix}{Line9}     \contestboxwhite{\myboxx}{Line10}
\contestbox{\myboxxi}{Line11}    \contestbox{\myboxxii}{Line12}
\contestbox{\myboxxiii}{Line13}

\newcommand{\TT}[1]{{\ttfamily #1}}
\newcommand{\BF}[1]{\textbf{#1}}
\newcommand{\INF}{\mathrm{INF}}
\lstset{basicstyle=\ttfamily\small,keywordstyle=\bfseries\color{deepblue},
  stringstyle=\color{deepgreen},numbers=left,
  numberstyle=\small\color{halfgray},frame=shadowbox}

\begin{document}
\begin{frame}
  \titlepage
\end{frame}

\begin{frame}{目录}
  \scriptsize
  \begin{multicols}{2}
    \tableofcontents[sectionstyle=show,subsectionstyle=hide]
  \end{multicols}
\end{frame}

\section{总体}
\subsection{难度与知识点}
\begin{frame}{难度预估}
  \begin{center}
  \large
  \begin{tabular}{cc}
    \toprule
    难度 & 题号 \\
    \midrule
    Easy & \myboxx{\BF{\TT{J}}}\ \myboxvi{\BF{\TT{F}}}\\
    EasyMedium & \myboxiii{\BF{\TT{C}}}\ \myboxix{\BF{\TT{I}}}\ \myboxxii{\BF{\TT{L}}}\\
    Medium & \myboxviii{\BF{\TT{H}}}\ \myboxiv{\BF{\TT{D}}}\ \myboxi{\BF{\TT{A}}}\\
    MediumHard & \myboxxiii{\BF{\TT{M}}}\ \myboxv{\BF{\TT{E}}}\ \myboxxi{\BF{\TT{K}}}\\
    Hard & \myboxii{\BF{\TT{B}}}\ \myboxvii{\BF{\TT{G}}}\\
    \bottomrule
  \end{tabular}
  \end{center}
\end{frame}

\begin{frame}{技巧表}
  \scriptsize
  \begin{center}
  \begin{tabular}{cccc}
    \toprule
    题号 & 核心技巧 & 题号 & 核心技巧\\
    \midrule
    \myboxi{\BF{\TT{A}}} & AC 自动机，矩阵快速幂 &
    \myboxviii{\BF{\TT{H}}} & 建图，BFS\\
    \myboxii{\BF{\TT{B}}} & 最大权闭合子图，最小割 &
    \myboxix{\BF{\TT{I}}} & 贪心\\
    \myboxiii{\BF{\TT{C}}} & 二分答案，线性DP &
    \myboxx{\BF{\TT{J}}} & 语法\\
    \myboxiv{\BF{\TT{D}}} & 构造 &
    \myboxxi{\BF{\TT{K}}} & 凸包，旋转卡壳，线段树\\
    \myboxv{\BF{\TT{E}}} & 莫比乌斯反演，线性筛 &
    \myboxxii{\BF{\TT{L}}} & 拓扑排序\\
    \myboxvi{\BF{\TT{F}}} & 区间并 &
    \myboxxiii{\BF{\TT{M}}} & 树链剖分\\
    \myboxvii{\BF{\TT{G}}} & 启发式算法 & &\\
    \bottomrule
  \end{tabular}
  \end{center}
\end{frame}

% A
\begin{frame}{题意概述}
  \begin{center}试题 \myboxi{\BF{\TT{A}}} \BF{纽约提喻法}\end{center}
  长度 $L\le 10^9$ 的小写字符串中，至多 $10$ 个位置被指定字符；
  给出总长度不超过 $100$ 的禁止串，求合法字符串数。
\end{frame}

{\setbeamercolor{frametitle}{fg=black,bg=Line1}
\section[A]{Problem A -- 纽约提喻法}
\subsection{A 纽约提喻法}
\begin{frame}{AC 自动机状态}
  对所有禁止串建立 AC 自动机。状态 $u$ 表示当前前缀的最长后缀，
  它也是某个模式串的前缀。
  \begin{itemize}
    \item 补全 $go(u,c)$ 转移。
    \item 坏状态满足
    \[
      bad[u]=terminal[u]\lor bad[fail[u]].
    \]
    \item 只保留安全状态以及安全状态之间的转移。
  \end{itemize}
  状态数 $S\le1+\sum|s_i|\le101$，可以直接使用 $S\times S$ 矩阵。
\end{frame}

\begin{frame}{分段矩阵转移}
  定义任意字符转移矩阵
  \[
    M_{u,v}=\#\{c\in\{\texttt{a},\ldots,\texttt{z}\}:go(u,c)=v,\ v\text{ 安全}\}.
  \]
  对强制字符 $c$，定义 $F_c$：
  \[
    (F_c)_{u,v}=
    \begin{cases}1,&go(u,c)=v,\ v\text{ 安全};\\0,&\text{其他情况}.\end{cases}
  \]
  将约束按位置排序。初始行向量仅根状态为 $1$，设 $last=0$。
  对每个 $(p,c)$：
  \[
    dp\gets dp\cdot M^{p-last-1}\cdot F_c,\qquad last\gets p.
  \]
  最后执行 $dp\gets dp\cdot M^{L-last}$，答案是所有安全状态计数之和。
\end{frame}

\begin{frame}{正确性与复杂度}
  AC 状态精确记录当前后缀是否已匹配到禁止串；删去所有进入坏状态的边，
  恰好排除出现过禁止串的字符串。

  \vspace{0.5em}
  $M^r_{u,v}$ 是从 $u$ 连续填 $r$ 个任意字符并安全到达 $v$ 的方案数，
  $F_c$ 固定当前位置，因此上述乘积逐段且不重不漏地计数全部合法串。

  \vspace{0.5em}
  预处理 $M^{2^j}$ 后，每段只用行向量乘幂矩阵：
  \[
    O(S^3\log L+(K+1)S^2\log L)
  \]
  时间，空间 $O(S^2\log L)$。
\end{frame}}

% B
\begin{frame}{题意概述}
  \begin{center}试题 \myboxii{\BF{\TT{B}}} \BF{AIMC 实验室}\end{center}
  三维网格每个点选择 $0/1$ 态。除了单点奖励外，
  距离平方为 $1,3$ 的邻域在满足全异态条件、距离平方为 $2$ 的邻域在满足全同态条件也有奖励。
\end{frame}

{\setbeamercolor{frametitle}{fg=white,bg=Line2}
\section[B]{Problem B -- AIMC 实验室}
\subsection{B AIMC 实验室}
\begin{frame}{奇偶翻转}
  令原状态为 $s_u$，并设
  \[
    p_u=(x_u+y_u+z_u)\bmod2,\qquad z_u=s_u\oplus p_u.
  \]
  距离平方为 $1,3$ 的点奇偶性不同，距离平方为 $2$ 的点奇偶性相同。
  所有题目中的条件都变为：
  \[
    Q_{u,d}=\{u\}\cup S_d(u)\text{ 中所有 }z\text{ 全为 }0
    \quad\text{或全为 }1.
  \]
  当 $p_u=0$ 时，$z_u=0$ 对应 $(A_1,A_3,A_5,A_7)$，
  $z_u=1$ 对应 $(A_2,A_4,A_6,A_8)$；$p_u=1$ 时两组交换。
\end{frame}

\begin{frame}{最小割建模}
  约定原格点在源侧当且仅当 $z=0$。令
  \[
    W=\sum_u\sum_{i=1}^{8}A_i(u),\qquad\INF=W+1.
  \]
  单点奖励 $w(0),w(1)$ 连
  \[
    S\to u\ (w(0)),\qquad u\to T\ (w(1)).
  \]
  对集合 $Q$ 的全为 $0$ 奖励 $w_0$，建立辅助点 $a_0$：
  \[
    S\to a_0\ (w_0),\qquad a_0\to v\ (\INF)\quad(v\in Q).
  \]
  对全为 $1$ 奖励 $w_1$，建立
  \[
    a_1\to T\ (w_1),\qquad v\to a_1\ (\INF)\quad(v\in Q).
  \]
\end{frame}

\begin{frame}{为何是总奖励减最小割}
  若 $Q$ 并非全在源侧，$a_0$ 不能留在源侧而切断无穷边，
  只能放到汇侧并割掉 $S\to a_0$，恰好失去 $w_0$。
  全为 $1$ 的情形对称。

  \vspace{0.5em}
  $\INF>W$ 保证最小割不会切无穷边。因此每个有限割的容量恰为
  未获得奖励之和，答案为
  \[
    \boxed{W-\operatorname{mincut}(S,T)}.
  \]

  \vspace{0.5em}
  每个格点的三类邻域各建两点，图有 $7V+2$ 个点、$O(V)$ 条边，
  $V=NMH\le8000$。容量和答案使用 \TT{long long}。
\end{frame}}

% C
\begin{frame}{题意概述}
  \begin{center}试题 \myboxiii{\BF{\TT{C}}} \BF{一分为二 2}\end{center}
  任选切分点 $k$，得到
  \[
    [a_1,\ldots,a_k,b_{k+1},\ldots,b_n].
  \]
  从中选恰好 $m$ 个互不相邻元素，最大化所选元素的最小值。
\end{frame}

{\setbeamercolor{frametitle}{fg=white,bg=Line3}
\section[C]{Problem C -- 一分为二 2}
\subsection{C 一分为二 2}
\begin{frame}{二分阈值与前后缀 DP}
  二分答案 $X$。只把值 $\ge X$ 的位置视为可选。
  对 $a$ 的前缀维护
  \[
  \begin{aligned}
  P_i^0&=\text{前 }i\text{ 个且 }i\text{ 不选的最大数量},\\
  P_i^1&=\text{前 }i\text{ 个且 }i\text{ 选的最大数量}.
  \end{aligned}
  \]
  \[
    P_i^0=\max(P_{i-1}^0,P_{i-1}^1),\qquad
    P_i^1=\begin{cases}P_{i-1}^0+1,&a_i\ge X;\\-\infty,&a_i<X.\end{cases}
  \]
  对 $b$ 的后缀同理维护 $Q_i^0,Q_i^1$。
\end{frame}

\begin{frame}{拼接边界}
  对固定切分 $k$，只有 $a_k,b_{k+1}$ 跨界相邻，故
  \[
  best(k)=\max\left\{
    P_k^0+\max(Q_{k+1}^0,Q_{k+1}^1),\
    P_k^1+Q_{k+1}^0
  \right\}.
  \]
  某个 $best(k)\ge m$ 当且仅当阈值 $X$ 可行。

  \vspace{0.5em}
  $P_0^0=Q_{n+1}^0=0$，另一状态为 $-\infty$，自然涵盖 $k=0,n$。
  判定具有单调性，整数二分约 $31$ 次：
  \[
    O(n\log10^9)
  \]
  时间，$O(n)$ 空间。
\end{frame}}

% D
\begin{frame}{题意概述}
  \begin{center}试题 \myboxiv{\BF{\TT{D}}} \BF{位距！？}\end{center}
  构造 $n\le5000$ 个平面点，点间距离至少 $10^{-4}$，
  并有至少
  \[
    \left\lfloor\tfrac12n\log_2n-n\right\rfloor
  \]
  对点的距离恰好为 $1$。
\end{frame}

{\setbeamercolor{frametitle}{fg=black,bg=Line4}
\section[D]{Problem D -- 位距！？}
\subsection{D 位距！？}
\begin{frame}{二进制超立方体投影}
  令 $K=\lceil\log_2n\rceil\le13$，$\varepsilon=2\times10^{-4}$，
  并定义单位向量
  \[
    \vec v_j=\left(\sqrt{1-(\varepsilon2^j)^2},\ \varepsilon2^j\right)
    \quad(0\le j<K).
  \]
  对 $i=0,1,\ldots,n-1$，输出
  \[
    P_i=\sum_{j=0}^{K-1}[\text{$i$ 的第 $j$ 位为 }1]\vec v_j.
  \]
  若仅第 $j$ 位不同，两点之差为 $\pm\vec v_j$，距离恰为 $1$。
  同时 $y(P_i)=\varepsilon i$，所以不同点的纵坐标差至少
  $2\times10^{-4}$。
\end{frame}

\begin{frame}{单位边计数}
  每个 $i$ 的每个 $1$ 位可以清零，得到更小的合法编号并产生一条单位边。
  因而单位边数至少
  \[
    F(n)=\sum_{i=0}^{n-1}\operatorname{popcount}(i).
  \]
  令 $n=2^q+r$，$0\le r<2^q$，有
  \[
    F(n)=q2^{q-1}+r+F(r).
  \]
  由此归纳得到
  \[
    F(n)\ge\tfrac12n\log_2n-n,
  \]
  正好达到要求。
  构造时间 $O(n\log n)$。
\end{frame}}

% E
\begin{frame}{题意概述}
  \begin{center}试题 \myboxv{\BF{\TT{E}}} \BF{根除}\end{center}
  定义 $f(n)$ 为 $n=ab^2$ 中最小的 $a$。这就是 $n$ 的平方自由核。
  查询
  \[
    \sum_{i=l}^{r}f(i),\qquad r\le10^{12},
  \]
  但至多 $10$ 个询问的右端点超过 $10^7$。
\end{frame}

{\setbeamercolor{frametitle}{fg=white,bg=Line5}
\section[E]{Problem E -- 根除}
\subsection{E 根除}
\begin{frame}{数论变形}
  令
  \[
    F(N)=\sum_{i=1}^{N}f(i),\qquad T(x)=\frac{x(x+1)}2.
  \]
  平方自由数示性函数为 $\mu^2$，故
  \[
    F(N)=\sum_{b\le\sqrt N}\sum_{a\le N/b^2}a\mu^2(a).
  \]
  利用
  \[
    \mu^2(a)=\sum_{d^2\mid a}\mu(d)
  \]
  得到
  \[
    \sum_{a\le X}a\mu^2(a)
    =\sum_{d\le\sqrt X}d^2\mu(d)
      T\left(\left\lfloor\frac{X}{d^2}\right\rfloor\right).
  \]
\end{frame}

\begin{frame}{合并平方因子}
  令 $m=bd$ 并交换求和：
  \[
    F(N)=\sum_{m\le\sqrt N}g(m)
      T\left(\left\lfloor\frac{N}{m^2}\right\rfloor\right),
    \qquad
    g(m)=\sum_{d\mid m}d^2\mu(d).
  \]
  $g$ 是积性函数，且
  \[
    g(m)=\prod_{p\mid m}(1-p^2).
  \]
  因此大端点只需枚举到 $\sqrt N\le10^6$。
  每个询问回答 $F(r)-F(l-1)$，在模 $998244353$ 下计算。
\end{frame}

\begin{frame}{预处理与复杂度}
  \begin{itemize}
    \item 对 $N\le10^7$，线性筛出 $f(i)$ 并做前缀和。
    \item 设 $p=\operatorname{lp}(i)$、$y=i/p$，则
    \[
      f(i)=
      \begin{cases}
        f(y)/p,&p\mid f(y);\\
        f(y)p,&p\nmid f(y).
      \end{cases}
    \]
    \item 对 $g$ 同样线性筛。若 $p\mid i$，$g(ip)=g(i)$；
    否则乘上 $1-p^2$。
  \end{itemize}
  预处理 $O(10^7)$。小端点 $O(1)$；大端点 $O(\sqrt N)$，
  至多约 $20$ 次，总枚举量不超过 $2\times10^7$。
\end{frame}}

% F
\begin{frame}{题意概述}
  \begin{center}试题 \myboxvi{\BF{\TT{F}}} \BF{凸包}\end{center}
  给出 $n$ 个严格一元二次不等式
  \[
    a_ix^2+b_ix+c_i>0,
  \]
  判断是否存在同时满足它们的整数 $x$。
\end{frame}

{\setbeamercolor{frametitle}{fg=black,bg=Line6}
\section[F]{Problem F -- 凸包}
\subsection{F 凸包}
\begin{frame}{每条不等式都是区间限制}
  记 $q(x)=ax^2+bx+c$。
  \begin{itemize}
    \item $a<0$ 时，$q(x)>0$ 的整数解是一个连续区间，或为空。
    \item $a>0$ 时，$q(x)\le0$ 的整数解是一个连续禁区，或为空。
  \end{itemize}
  先相交所有 $a<0$ 的可行区间，得到总候选区间 $[L,R]$；
  再把所有 $a>0$ 的禁区合并，判断它们是否覆盖 $[L,R]$。

  \vspace{0.5em}
  若不存在 $a<0$ 的式子，直接输出 \TT{YES}：
  取足够大的 $x$ 时所有开口向上的二次式均为正。

  \vspace{0.5em}
  严格不等号很重要：开口向上的式子必须删去 $q(x)\le0$，
  根处取等同样不合法。

  \vspace{0.5em}
  时间$O(n\log n)$，空间 $O(n)$。
\end{frame}}

% G
\begin{frame}{题意概述}
  \begin{center}试题 \myboxvii{\BF{\TT{G}}} \BF{星之彼端}\end{center}
  通信题。第一次运行选出 $200000$ 个诱导四元环；
  第二次收到其中一个环删去任意一点后的三元组，需恢复环编号。

  \vspace{1em}
  图公开且固定，允许针对给定随机图做确定性的构造。
\end{frame}

{\setbeamercolor{frametitle}{fg=black,bg=Line7}
\section[G]{Problem G -- 星之彼端}
\subsection{G 星之彼端}
\begin{frame}{唯一解码条件}
  若任意两个输出的四元环至多共享两个顶点，则它们不会共享一个三元组。
  对每个选中环的四个三元组排序后建立哈希表
  \[
    \{\text{三元组}\}\longrightarrow\{\text{输出编号}\},
  \]
  第二次运行即可唯一回答。

  \vspace{0.5em}
  因而任务变成：从所有诱导 $C_4$ 中选 $200000$ 个，
  且四个三元组均不与已选环冲突。
\end{frame}

\begin{frame}{枚举所有诱导四元环}
  设 $u,v$ 是一对不相邻的对角点，枚举其公共邻居 $x,y$。
  若 $x,y$ 也不相邻，则
  \[
    u-x-v-y-u
  \]
  是诱导四元环。邻接矩阵或 bitset 可 $O(1)$ 判边。

  \vspace{0.5em}
  实现中可枚举长度为 $2$ 的路径，把共同中点按端点对分组；
  对四元组排序并按规范规则去重。令
  \[
    W=\sum_v\binom{\deg v}{2},\qquad
    K=\sum_{u<v}\binom{|N(u)\cap N(v)|}{2},
  \]
  枚举量约为 $O(W+K)$。
\end{frame}

\begin{frame}{三元组打包}
  对每个候选环记录四个三元组，统计每个三元组出现次数。
  定义环的冲突度为其四个三元组频次之和，按
  冲突度和规范化四元组的字典序贪心：
  \begin{itemize}
    \item 四个三元组都未占用时选入。
    \item 对未选环，若它所有冲突三元组只归属于同一个已选环，可替换该环。
    \item 再执行一轮可直接加入的贪心，必要时重复。
  \end{itemize}
  这是公开固定数据上的构造性启发式，不是任意图上的通用保证。
\end{frame}}

% H
\begin{frame}{题意概述}
  \begin{center}试题 \myboxviii{\BF{\TT{H}}} \BF{一般攻击魔法}\end{center}
  空格可填正斜杠镜、反斜杠镜或保持空，固定格不能改。
  判断能否令光线从给定边界端口射入后，恰好从目标端口射出并输出方案。
\end{frame}

{\setbeamercolor{frametitle}{fg=white,bg=Line8}
\section[H]{Problem H -- 一般攻击魔法}
\subsection{H 一般攻击魔法}
\begin{frame}{把每个格子拆成四个端口}
  每个格子拆为 $U,D,L,R$ 四个状态；相邻格子对应端口相连：
  \[
    R_{i,j}\leftrightarrow L_{i,j+1},\qquad
    D_{i,j}\leftrightarrow U_{i+1,j}.
  \]
  固定格内部只连它实际的两对端口：
  \[
  \begin{array}{c|c}
    \text{填法}&\text{内部配对}\\ \hline
    \texttt{.}&U-D,\ L-R\\
    \texttt{/}&U-L,\ D-R\\
    \text{反斜杠}&U-R,\ D-L
  \end{array}
  \]
  原本为空的格子暂时把四个端口连成 $K_4$。在此松弛图中从起点 BFS 到终点。
\end{frame}

\begin{frame}{从最短路恢复填法}
  恢复 BFS 的最短简单路径。若路径在空格中使用端口对：
  \[
  \begin{array}{c|c}
    U-D\text{ 或 }L-R&\texttt{.}\\
    U-L\text{ 或 }D-R&\texttt{/}\\
    U-R\text{ 或 }D-L&\text{反斜杠}
  \end{array}
  \]
  就填入对应状态；未经过的空格统一填 \TT{.}。

  \vspace{0.5em}
  若最短路在同一空格的 $K_4$ 内使用两次，取第一次和最后一次经过的端口，
  $K_4$ 中有直接边，可以严格缩短路径，矛盾。因此每个可变格至多需要
  一条内部配对边。
\end{frame}

\begin{frame}{正确性与复杂度}
  任意真实填法的每条格内连边都是 $K_4$ 的边，故真实可行必能在松弛图中找到。
  反过来，恢复表保留了最短路所用的端口对；每个端口只有唯一的格内配对，
  光线不会在路径中途偏离，因而确实能到达目标。

  \vspace{0.5em}
  图有 $4nm$ 个点、$O(nm)$ 条边，BFS、恢复和输出均为
  \[
    O(nm).
  \]
\end{frame}}

% I
\begin{frame}{题意概述}
  \begin{center}试题 \myboxix{\BF{\TT{I}}} \BF{天外来物}\end{center}
  依任意顺序采集一条链上的水晶。采集一个点会翻转尚未采集的相邻点能量，
  求最大总能量。
\end{frame}

{\setbeamercolor{frametitle}{fg=black,bg=Line9}
\section[I]{Problem I -- 天外来物}
\subsection{I 天外来物}
\begin{frame}{所有符号模式只受一个奇偶条件限制}
  记水晶 $i$ 被采集时的符号为 $(-1)^{x_i}$。
  每条相邻边恰好在较晚被采集的一个端点处贡献一次翻转，因此
  \[
    x_1\oplus x_2\oplus\cdots\oplus x_n=(n-1)\bmod2.
  \]

  \vspace{0.5em}
  反过来，任意满足此奇偶条件的符号模式都能实现：将每条边按采集先后定向，
  指定每个点入边数的奇偶性为 $x_i$。在路径上从左到右即可唯一递推边方向；
  总奇偶条件正是最后一条边一致的条件。路径任意定向无环，取拓扑序就是一种采集顺序。
\end{frame}

\begin{frame}{贪心取最优符号}
  想让每项贡献为 $|a_i|$：正数取正号，负数取负号。
  设原数组负数个数为 $cnt$。
  \begin{itemize}
    \item 若 $cnt\bmod2=(n-1)\bmod2$，全部绝对值都能获得。
    \item 否则必须翻转奇数个目标符号，最优只翻转绝对值最小的一项。
  \end{itemize}
  \[
    \boxed{\sum_i|a_i|-
    \begin{cases}
      0,&cnt\bmod2=(n-1)\bmod2;\\
      2\min_i|a_i|,&\text{否则}.
    \end{cases}}
  \]
  每组 $O(n)$ 时间、$O(1)$ 额外空间。
\end{frame}}

% J
\begin{frame}{题意概述}
  \begin{center}试题 \myboxx{\BF{\TT{J}}} \BF{时间谜题 2}\end{center}
  将非负整数时间 $A{:}B$ 规范化为唯一的
  $0\le h<24$、$0\le m<60$ 的标准时间。
\end{frame}

{\setbeamercolor{frametitle}{fg=white,bg=Line10}
\section[J]{Problem J -- 时间谜题 2}
\subsection{J 时间谜题 2}
\begin{frame}{直接按总分钟数取模}
  两种等价操作都不改变总分钟数
  \[
    60A+B
  \]
  模 $24\times60=1440$ 的值。因此
  \[
    t=(60A+B)\bmod1440,\qquad
    h=\left\lfloor\frac{t}{60}\right\rfloor,\quad m=t\bmod60.
  \]
  也可写作
  \[
    h=(A+\lfloor B/60\rfloor)\bmod24,\qquad m=B\bmod60.
  \]

  \vspace{1em}
  每组只需 $O(1)$。
\end{frame}}

% K
\begin{frame}{题意概述}
  \begin{center}试题 \myboxxi{\BF{\TT{K}}} \BF{永世沙漏}\end{center}
  给定逆时针简单多边形，统计稳定方向 $d$。
  需要找到一条垂直于 $d$ 的线 $l$，
  使多边形中每一点都能沿 $d$ 方向直达 $l$ 而不离开多边形。
\end{frame}

{\setbeamercolor{frametitle}{fg=black,bg=Line11}
\section[K]{Problem K -- 永世沙漏}
\subsection{K 永世沙漏}
\begin{frame}{投影坐标下的充要条件}
  对候选整数方向 $\mathbf d=(a,b)$，令
  \[
    \mathbf e=(b,-a),\qquad X(P)=P\cdot\mathbf e,\qquad Y(P)=P\cdot\mathbf d.
  \]
  直线 $l$ 形如 $Y=c$，投影线平行 $\mathbf d$。

  \vspace{0.5em}
  稳定当且仅当：
  \begin{enumerate}
    \item 多边形关于 $X$ 单调，即每条 $X=\mathrm{const}$ 截线是线段；
    \item 所有截线有公共点 $(X,c)$。
  \end{enumerate}
  若两条边界链为下链 $D$ 和上链 $U$，第二项等价于
  \[
    \max_{P\in D}Y(P)\le\min_{P\in U}Y(P).
  \]
\end{frame}

\begin{frame}{候选方向只有线性个}
  考虑凸包。稳定直线在 $X$ 的两个极值处接触凸包：
  \begin{itemize}
    \item 若某个接触面是一条边，则该边平行于 $\mathbf d$。
    \item 若两端都是顶点，则这对顶点是凸包的一对对踵点，且
    $\mathbf d$ 垂直于它们的连线。
  \end{itemize}
  因而只需枚举
  \[
    \text{凸包边方向}\quad\text{和}\quad
    \text{凸包对踵点连线的垂直方向}.
  \]
  旋转卡壳在线性时间列出全部对踵关系。规范化、去重，并只保留
  $a>0,b\ge0$ 的方向，候选数为 $O(n)$。
\end{frame}

\begin{frame}{快速验证一个候选}
  对原多边形的循环顶点序列建立两棵线段树：
  \begin{itemize}
    \item 顶点树：节点维护顶点集合凸包，查询区间内 $P\cdot\mathbf d$、
    $P\cdot\mathbf e$ 的最值；
    \item 边树：节点维护边向量集合凸包，查询区间内
    $\Delta P\cdot\mathbf e$ 的最值。
  \end{itemize}
  找到 $X$ 极值支撑面后，按原边界验证四段：
  \[
  L^-\to R^-:\Delta X\ge0,\quad R^-\to R^+:\Delta X=0,
  \]
  \[
  R^+\to L^+:\Delta X\le0,\quad L^+\to L^-:\Delta X=0,
  \]
  最后检查下链的 $\max Y$ 不超过上链的 $\min Y$。
\end{frame}

\begin{frame}{复杂度与精度}
  一个环形范围拆成至多两个普通区间，线段树访问 $O(\log n)$ 个节点；
  每个局部凸包上二分支持函数极值，故一次验证 $O(\log^2 n)$。

  \vspace{0.5em}
  凸包、旋转卡壳与全部候选总计 $O(n)$，建树 $O(n\log n)$，总复杂度
  \[
    O(n\log^2 n).
  \]
  点和方向均使用整数表示，点积、叉积使用 \TT{\_\_int128}，避免浮点误差。
\end{frame}}

% L
\begin{frame}{题意概述}
  \begin{center}试题 \myboxxii{\BF{\TT{L}}} \BF{净化行动 3}\end{center}
  无限棋盘上有若干炮。若同一行或列中两个炮之间恰好隔一个炮，
  它们互相攻击。每次只能删除当前未受攻击的炮，判断能否删空。
\end{frame}

{\setbeamercolor{frametitle}{fg=black,bg=Line12}
\section[L]{Problem L -- 净化行动 3}
\subsection{L 净化行动 3}
\begin{frame}{同一行或列的结构}
  将同一横线或纵线上的炮按坐标排序。
  \begin{itemize}
    \item 若初始有至少 $4$ 个炮，则该线上的每个炮都与相隔一个的炮互相攻击，
    它们一个也无法删除，答案必为 \TT{NO}。
    \item 若恰有 $3$ 个炮，只有两端受攻击；中间炮必须先于两端删除。
    \item 若至多 $2$ 个炮，没有攻击关系。
  \end{itemize}
  删除不会制造新的隔一个炮的攻击关系；只会消灭已有关系。
\end{frame}

\begin{frame}{拓扑序即删除顺序}
  对每条恰有三个炮的横线或纵线，向中间炮到两个端点各连一条有向边：
  \[
    middle\longrightarrow endpoint.
  \]
  所有约束同时满足当且仅当有向图无环。

  \vspace{0.5em}
  任意拓扑序中，一个点被删前，所有要求它先删除的中间炮都已消失，
  因而它当前不受攻击；反之任何合法删除顺序也必须遵守这些边。

  \vspace{0.5em}
  按 $x$、按 $y$ 分组排序后建边，再拓扑排序：
  \[
    O(n\log n)
  \]
  时间、$O(n)$ 空间。
\end{frame}}

% M
\begin{frame}{题意概述}
  \begin{center}试题 \myboxxiii{\BF{\TT{M}}} \BF{魔女审判}\end{center}
  树上每个点有一个当前唯一的权值。支持单点换权；
  查询值域 $[l,r]$ 中所有权值所在节点，是否恰好组成两条向上路径的并。
\end{frame}

{\setbeamercolor{frametitle}{fg=black,bg=Line13}
\section[M]{Problem M -- 魔女审判}
\subsection{M 魔女审判}
\begin{frame}{查询先确定关键端点}
  离线收集初值和所有更新值并离散化，维护一棵按权值建的线段树。
  每个节点记录
  \[
    (cnt,\ \sum dep,\ \text{最浅深度及其个数和节点},\ \text{一个最深节点}).
  \]
  查询 $[l,r]$ 前先检查
  \[
    cnt=r-l+1.
  \]
  否则某个整数权值不存在，立刻为 \TT{NO}。最浅点还必须唯一；
  记它为 $z$，取一个最深点为 $x$，并检查 $z$ 是 $x$ 的祖先。
\end{frame}

\begin{frame}{从深度和恢复另一条链}
  令
  \[
    v=r-l+1,\quad len_1=dep_x-dep_z+1,\quad len_2=v-len_1.
  \]
  第一条 $z\to x$ 的深度和为
  \[
    S_1=dep_z\,len_1+\frac{len_1(len_1-1)}2.
  \]
  先用树链剖分确认整条 $z\to x$ 的权值都在 $[l,r]$。
  若 $len_2=0$，直接为 \TT{YES}。

  \vspace{0.5em}
  否则另一端点 $y$ 的深度唯一：
  \[
    dep_y=\frac{\sum dep-S_1+len_2(len_2-1)/2}{len_2}.
  \]
  要求整除且 $dep_z+len_2\le dep_y\le dep_x$。
\end{frame}

\begin{frame}{定位 $y$ 并验证}
  对每个深度维护当前权值的有序集合。主链 $z\to x$ 在深度 $dep_y$
  有唯一节点 $p$。合法时，$[l,r]$ 内该深度恰有两个节点：$p$ 与 $y$；
  用 \TT{lower\_bound} 至多检查三个元素即可判断并找出 $y$。

  \vspace{0.5em}
  令 $w=\operatorname{LCA}(x,y)$，还需检查
  \[
    z\text{ 是 }y\text{ 的祖先},\qquad dep_w=dep_y-len_2.
  \]
  再验证 $w\to y$ 的每个点权都在 $[l,r]$。两条路径的并有
  $len_1+len_2=v$ 个不同点，且全部落在值域内；由 $cnt=v$ 和权值唯一性，
  该并集恰为查询集合。
\end{frame}

\begin{frame}{数据结构与复杂度}
  同时维护：
  \begin{itemize}
    \item 值域线段树：计数、深度和、最浅和最深节点；
    \item 树链剖分上的线段树：路径权值最小值和最大值；
    \item 每个深度一个 \TT{set}：定位给定值域内的候选节点；
    \item 倍增 LCA。
  \end{itemize}
  更新时同步修改三处。值域查询和更新为 $O(\log n)$，
  两次树链查询为 $O(\log^2 n)$，总复杂度
  \[
    O((n+q)\log^2 n).
  \]
\end{frame}}

\section*{End}
\begin{frame}{Thanks}
  \begin{center}\Huge Thanks!\end{center}
\end{frame}
\end{document}
